We volgen daarbij de grote structuur van de metadatafiche (nog voor discussie vatbaar):
1.	Context (gebruikersbehoeften en conceptualisering)
  - falsificatie? Niet enkel data om beleid te ondersteunen maar ook voor oppositie -> Reinhard geeft insteek
  - belangrijk is om goede dataontsluiting te doen, niet enkel statistiek
2.	Methoden (operationalisering)
  - klemtoon op operationeel doel van admin data
  - primair gebruik van data
  - vb. EPC data -> operationeel administratie
  		werkloosheidsdata -> wordt geherdefinieerd
  
3.	Kwaliteit
4.	Ontsluiting
5.	Administratieve gegevens & algemeen beleid




Ter voorbereiding kunnen jullie al opschrijven wat volgens jullie zeker ergens moet worden uitgelegd en ook interessante voorbeelden noteren (bv. suggestie Reinhard om te verduidelijken dat ontsluiting pure administratieve data niet automatisch leidt tot kwaliteitsvolle statistieken). 
Het meest uitdagend zullen punten 2 en 3 zijn, punt 1 moeten we nog feedback Jo integreren, puten 4 & 5 lijken me redelijk rechttoe rechtaan.

Ik heb zelf al wat suggesties:
-	Vanuit klantvriendelijkheid hulp bij interpretatie en inhoud zo veel mogelijk naar boven (context > methoden > kwaliteit), pure administratieve info zo veel mogelijk naar beneden (hoe ontsloten, admin gegevens, en algemeen beleidskader rond productieproes)
-	Verder durven loslaten van SIMS-structuur. Ik loop bv. nog vast op de methodesectie, waar variabelen worden besproken in één blok en bewerkingen in een ander. Dat verhindert een vloeiend geschreven uitleg over het productieproces. We moeten de SIMS vooral zien als een check list, niet als een structureel raamwerk.
